% seb.tex -- Smallest enclosing balls: a companion note to cvxball.
% Build:  pdflatex seb.tex   (twice)   or   tectonic seb.tex

\documentclass[10pt,a4paper]{article}

\usepackage[T1]{fontenc}
\usepackage{amsmath,amssymb,amsthm}
\usepackage{mathtools}
\usepackage{booktabs}
\usepackage{xcolor}
\usepackage[margin=2.1cm]{geometry}
\usepackage[colorlinks=true,allcolors=blue!50!black]{hyperref}

\theoremstyle{plain}
\newtheorem{proposition}{Proposition}

\DeclareMathOperator{\conv}{conv}
\DeclareMathOperator{\aff}{aff}
\DeclareMathOperator*{\argmax}{arg\,max}
\newcommand{\R}{\mathbb{R}}
\newcommand{\norm}[1]{\lVert #1 \rVert}
\newcommand{\T}{^{\mathsf{T}}}

\setlength{\parskip}{0pt}
\title{\vspace{-1.2cm}\bfseries Two active-set methods for smallest enclosing balls}
\author{Thomas Schmelzer}
\date{\today}

\begin{document}
\maketitle
\vspace{-0.6cm}

\begin{abstract}
\noindent
Two active-set methods solve the smallest enclosing ball problem from opposite
sides: the dual ascent over the unit simplex of Dearing and
Zeck~\cite{dearing2009}, whose ball encloses nothing until it terminates, and the
primal deflation of Fischer, G\"artner and Kutz~\cite{fischer2003}, which is
feasible at every step. This note compares the two with their linear algebra held
fixed. They keep the same support set, solve the same subproblem, and can repair
the same $QR$ factorisation of the same edge matrix in $O(dr)$; handing both that
one data structure removes the usual confound, so what separates them is a
property of the two searches rather than of two implementations. On Gaussian
clouds to $d=1.6\cdot10^4$ they agree on the ball and on its support on every row,
and the primal method takes $1.2$ to $2.1$ times as many steps --- the price of
never taking a full step --- which the wall clock damps to a factor of $1.1$ to
$1.6$. The cone program both replace, handed to an interior-point solver, is two
to three orders of magnitude slower where it runs at all, and stops at a tolerance
where both of these terminate exactly. Neither method is ours; both ship in
\texttt{cvxball}.
\end{abstract}

\section{Introduction}
\label{sec:intro}

Sylvester stated the problem in 1857 in a single sentence~\cite{sylvester1857}:
given a finite set of points in the plane, find the smallest circle containing
all of them. It is the $1$-centre facility location problem; it is the primitive
under bounding-sphere hierarchies in graphics and collision detection; and as the
minimum enclosing ball of a set of feature vectors it is the core-set problem
behind kernel machines.

The algorithmic history runs in two lines that have largely stayed apart. The
first is computational-geometric. Welzl's randomised incremental
method~\cite{welzl1991} solves the problem in expected $O(n)$ time for fixed $d$,
with a constant that grows superexponentially in $d$; G\"artner~\cite{gaertner1999}
made its predicates robust in floating point, G\"artner and
Sch\"onherr~\cite{gaertner2000} recast the problem as a quadratic program and
solved it exactly, and Fischer, G\"artner and Kutz~\cite{fischer2003} left the
recursion behind for a pivoting method that deflates an enclosing ball, carrying a
dynamic $QR$ of the support's edge matrix and reporting $d$ in the thousands --- a
finite algorithm of that second kind going back to Elzinga and
Hearn~\cite{elzinga1972}. The second line is an optimization one, reaching the same
object from the dual: Dearing and Zeck~\cite{dearing2009} gave a dual algorithm
maintaining a support set and raising the radius to $R$ from below, and Cavaleiro
and Alizadeh brought its iteration from $O(d^3)$ to $O(d^2)$ by updating a
factorisation rather than rebuilding it~\cite{cavaleiro2018} and then extended the
method to balls~\cite{cavaleiro2021}.

This note states the two methods in a common notation
(Sections~\ref{sec:dual}--\ref{sec:fgk}), shows that they can share one data
structure (Section~\ref{sec:qr}), and reports what separates them once that data
structure is held fixed (Section~\ref{sec:numerics}). Handing both the same linear
algebra removes the usual confound: what is left is a property of the two searches
rather than of two implementations. Neither method is new and none is claimed
here; what this note adds is the controlled comparison, and an implementation of
both that makes it reproducible.

\section{The problem and its two sides}
\label{sec:dual}

For $S=\{p_1,\dots,p_n\}\subset\R^d$ the smallest enclosing ball is the solution of
\begin{equation}
  \label{eq:primal}
  \min_{x,r}\ r \quad\text{s.t.}\quad \norm{p_i-x}\le r,\ i=1,\dots,n,
\end{equation}
a second-order cone program in $(r,x)\in\R^{1+d}$ with $n$ blocks
$(r,\,p_i-x)\in\mathcal{Q}^{d+1}$. Existence follows from coercivity of
$f(x)=\max_i\norm{p_i-x}$ and uniqueness from strict convexity of $f^2$; we write
$(R,x^\star)$ for the optimum.

Let $\Delta_n=\{u\in\R^n: u\ge0,\ \mathbf{1}\T u=1\}$, and for $u\in\Delta_n$ put
$x(u)=\sum_i u_ip_i$ and
\begin{equation}
  \label{eq:phi}
  \varphi(u)=\sum_i u_i\norm{p_i}^2-\norm{x(u)}^2 ,
\end{equation}
a concave quadratic. Everything below rests on the parallel-axis identity, valid
for all $y\in\R^d$ and $u\in\Delta_n$:
\begin{equation}
  \label{eq:id}
  \sum_i u_i\norm{p_i-y}^2
    \;=\; \varphi(u)+\norm{x(u)-y}^2 ,
\end{equation}
obtained by expanding and collecting, using $\sum_iu_ip_i=x(u)$ on the cross term
and $\mathbf{1}\T u=1$ on $\norm{y}^2$. Measuring from anywhere other than the
weighted mean costs precisely $\norm{x(u)-y}^2$, and $\varphi(u)$ is what is left
when $y$ is that mean.

\begin{proposition}
  \label{prop:duality}
  $R^2=\max_{u\in\Delta_n}\varphi(u)$, and $x^\star=x(w)$ for any maximiser $w$.
\end{proposition}

\noindent
Both directions are short: $\varphi\le R^2$ is~\eqref{eq:id} applied to any
enclosing ball, and equality follows by separating $x^\star$ from $\conv T$,
$T=\{p_i:\norm{p_i-x^\star}=R\}$, if the two were disjoint. See
also~\cite{gaertner2000}. Since
$\partial\varphi/\partial u_i=\norm{p_i-x(u)}^2-\norm{x(u)}^2$, the first-order
conditions read, after cancelling the common term,
\begin{equation}
  \label{eq:kkt}
  \norm{p_i-x}=R\ \ (w_i>0),\qquad \norm{p_i-x}\le R\ \ (w_i=0),
\end{equation}
a geometric certificate: the points carrying weight lie on the boundary sphere and
the rest are enclosed, which a caller checks in one pass. By Carath\'eodory the
support may be taken of size at most $d+1$, so~\eqref{eq:primal} is determined by
at most $d+1$ of the $n$ points and the problem is combinatorial in the choice of
that subset.

Read forwards,~\eqref{eq:kkt} is a search for weights: keep a set of points on the
sphere, adjust the weights, stop when no point lies outside. Read backwards it is
a search for balls: keep a ball containing $S$ with a set $T$ on its boundary, and
shrink it until its centre falls inside $\conv T$. The second rests on a fact due
to Seidel, recorded in~\cite{gaertner2000} and Lemma~1 of~\cite{fischer2003}: if
$T$ lies on the boundary of a ball $B$ centred at $c$, then $B$ is the smallest
ball enclosing $T$ if and only if $c\in\conv T$. With $S\subseteq B$ and
$T\subseteq S$ the inclusions $\mathrm{seb}(T)\subseteq\mathrm{seb}(S)\subseteq B$
close, so $\conv T$ is the only thing left to check.

Sections~\ref{sec:method} and~\ref{sec:fgk} are those two readings made into
methods. They keep the same kind of state --- a subset of $S$ that will end up on
the boundary, and a point in its affine hull --- and differ in which feasibility
they refuse to give up. The first stays dual feasible, so by
Proposition~\ref{prop:duality} its radius rises to $R$ from below and its ball
encloses nothing until the last iteration; the second stays primal feasible, so
its radius falls to $R$ from above and every iterate is an answer, merely not yet
the best one.

\section{The dual active-set method}
\label{sec:method}

This is the dual algorithm of Dearing and Zeck~\cite{dearing2009} as sharpened by
Cavaleiro and Alizadeh~\cite{cavaleiro2018,cavaleiro2021}, restated so that the
correspondence with Section~\ref{sec:fgk} is exact. Two details below depart from
that line of work and are given as differences rather than as claims: the step
rule here is a full step to the circumcentre of the current face followed by a
ratio test~\eqref{eq:ratio}, where~\cite{cavaleiro2021} performs an exact curve
search; and an affinely dependent support is resolved by a descent direction in
the null space rather than by dropping a point before the search begins.

Maintain a working set $F$ and weights $u\in\Delta_n$ supported on $F$ with
$\varphi(u)$ non-decreasing. This is a primal active-set method run on the
dual~\eqref{eq:phi}: the free variables are the weights of the support, every
subproblem solves one face exactly, and feasibility for~\eqref{eq:primal} arrives
only at the end.

\paragraph{Subproblem.} Maximising $\varphi$ over weights supported on $F$ with no
sign restriction is, by~\eqref{eq:kkt} without the inequalities, the requirement
that $x$ be equidistant from $\{q_j\}_{j\in F}$ and lie in their affine hull.
With $D$ the $(m-1)\times d$ matrix of edges $e_j=q_j-q_0$ and $x=q_0+D\T y$,
\begin{equation}
  \label{eq:sub}
  (DD\T)\,y=\tfrac12\bigl(\norm{e_1}^2,\dots,\norm{e_{m-1}}^2\bigr)\T,
  \qquad w=\Bigl(1-\textstyle\sum_j y_j,\ y\Bigr),
\end{equation}
a dense system of order $m-1\le d$, non-singular exactly when $F$ is affinely
independent. This is the circumcentre of the face in barycentric coordinates.

\paragraph{Dropping.} If $\min_j w_j<0$ the circumcentre lies outside
$\conv\{q_j\}$. Move along $s=w-u$ by the ratio test
\begin{equation}
  \label{eq:ratio}
  \alpha=\min\{-u_i/s_i:\ s_i<0\},
\end{equation}
which drives at least one weight to zero; remove it from $F$.

\paragraph{Degeneracy.} If $F$ is affinely dependent, $DD\T$ is singular. The
directions $s$ with $\mathbf{1}\T s=0$ and $D\T s_{1:}=0$ leave $x(u)$ fixed, so
$\varphi(u+ts)=\varphi(u)+t\sum_is_i\norm{p_i}^2$ is linear and the subproblem is
unbounded. Project the gradient onto that null space and apply~\eqref{eq:ratio},
shrinking $F$ by one.

\paragraph{Adding.} Otherwise accept $u\leftarrow w$, set $x=x(u)$ and
$R=\max_{i\in F}\norm{p_i-x}$, and let $k=\argmax_i\norm{p_i-x}$. If
$\norm{p_k-x}\le R$ then~\eqref{eq:kkt} holds and $(R,x,u)$ is optimal; otherwise
append $k$ to $F$ with weight zero. Since
$\partial\varphi/\partial u_k-\partial\varphi/\partial u_i
=\norm{p_k-x}^2-\norm{p_i-x}^2$, the farthest violator is the fixed variable with
the largest reduced gradient, so this is also the steepest-ascent choice.

Each iteration either strictly increases $\varphi$ or strictly shrinks $F$, and
there are finitely many working sets, so the method is finite; a degenerate
zero-length step needs an anti-cycling rule in exact arithmetic, and in practice an
iteration cap serves as the safety net.

\section{The primal pivoting method}
\label{sec:fgk}

The method of Fischer, G\"artner and Kutz~\cite{fischer2003} carries a ball rather
than a set of weights, and never surrenders it. Its state is a pair $(c,T)$, a
centre and a support set, subject to the invariant of their Fig.~2:
\begin{equation}
  \label{eq:fgk-invariant}
  S\subseteq B(c,T),\qquad T\subseteq\partial B(c,T),\qquad
  T\ \text{affinely independent},
\end{equation}
where $B(c,T)$ is the ball about $c$ through the farthest point of $T$. By
Seidel's lemma the pair is optimal as soon as $c\in\conv T$. Each iteration is a
dropping phase, entered only once $c$ has landed in $\aff(T)$, followed by a
walking phase.

\paragraph{Dropping.} $c\in\aff(T)\setminus\conv T$ means some affine coefficient
of $c$ with respect to $T$ is negative; drop such a point. The coefficients are
the same object as the barycentric weights of~\eqref{eq:sub} and come from the
same factorisation (Section~\ref{sec:qr}), but here they are read rather than
moved: there is no ratio test, because $c$ does not travel during a drop. Their
Lemmata~4 and~5 together say the dropped point cannot stop the walk that
immediately follows.

\paragraph{Walking.} Move $c$ along $e=\mathrm{cc}(T)-c$. Their Lemma~3 says that
segment is orthogonal to $\aff(T)$, that $T$ stays on the boundary the whole way,
and that the radius strictly decreases. Two consequences make the step cheap.
Orthogonality makes $\mathrm{cc}(T)$ the projection of $c$ onto $\aff(T)$, so no
circumsphere is ever solved for and one projection does the work
of~\eqref{eq:sub}. And the stopping fraction has a closed form: with $c(a)=c+ae$,
a point $p$ joins the boundary when $\norm{p-c(a)}$ equals the common support
distance; the $a^2$ terms cancel, and $\langle q-c,e\rangle=\norm{e}^2$ for every
$q\in T$ leaves
\begin{equation}
  \label{eq:fgk-step}
  a_p=\frac{R_c^2-\norm{p-c}^2}
           {2\bigl(\norm{e}^2-\langle p-c,\,e\rangle\bigr)},
  \qquad R_c=\max_{q\in T}\norm{q-c},
  \qquad a=\min\{1,\ \min_p a_p\}.
\end{equation}
Criterion~(1) of Lemma~4 --- that $p$ lies behind $\aff(T)$ and so cannot stop the
walk --- is exactly the statement that this denominator is non-positive, so one
sign test separates the candidates. If $a=1$ the centre arrives at
$\mathrm{cc}(T)$, lands in $\aff(T)$, and the next iteration drops; otherwise the
minimising point enters $T$ and the ball has shrunk.

\paragraph{Pivot rules.} Ties are resolved two ways. Bland's rule, adapted --- fix
an order on $S$, take the smallest-rank candidate at both the drop and the
admission --- is what their Theorem~1 proves terminating, degeneracies included,
and is also recorded there as slow. The heuristic, which their own code runs and
which Table~\ref{tab:highd} reports, drops the \emph{most} negative coefficient and
admits the tied point farthest from $\aff(T)$, leaving the support best
conditioned. The radius falls strictly on every walk and no configuration recurs,
so the method is finite.

\paragraph{How the two differ.} The support sets agree and the drop criteria agree
--- a negative barycentric coordinate, in both --- and the mechanics agree nowhere
else. Section~\ref{sec:method} takes the whole step to $\mathrm{cc}(F)$ and only
then consults the data, admitting the farthest violator; this method never takes a
full step while any point is still outside, truncating the walk at whichever point
first touches the shrinking sphere, which is also the point that enters. A drop
here is immediate and followed by a resumed walk; there it is a ratio test in
weight space, moving every weight at once. The invariant~\eqref{eq:fgk-invariant}
also requires $T$ affinely independent, with no fallback if it is not, so an
entering point must be certified off $\aff(T)$ before admission;
Section~\ref{sec:method} instead permits a dependent support and answers it with a
null-space direction. Neither choice dominates --- one avoids the case, the other
resolves it.

\section{The shared factorisation}
\label{sec:qr}

Per iteration, both methods want the same two things from their support: whether
it is affinely independent, and where the circumcentre of its face lies. Both read
those off the same object. Let $D$ be the $r\times d$ matrix of edges $q_j-q_0$ of
the current support and let $D\T=QR$ be its economic factorisation. Three
identities take every per-iteration quantity off $d$ and onto the $r\times r$
block:
\begin{equation}
  \label{eq:qr-identities}
  DD\T=R\T R,\qquad \ker D\T=\ker R,\qquad \norm{q_j-q_0}^2=\norm{R e_j}^2 .
\end{equation}
The first makes~\eqref{eq:sub} two triangular solves against a Cholesky factor
already to hand, with no $O(r^2d)$ product; the second makes the
affine-dependence test a rank test on $R$ --- note it is the \emph{right} null
space, not the one the edge matrix itself would hand you; the third supplies the
right-hand side. The pivoting method reads the same $R$ for its projection and its
coefficients.

Rebuilding that factorisation costs $O(dr^2)$ and repairing it costs $O(dr)$. The
repair is section~4 of~\cite{fischer2003}, and the same move on the dual side
is~\cite{cavaleiro2018}. Three updates cover every move either method makes: a
point joining the support appends a column, a point other than $q_0$ leaving
deletes one, and $q_0$ leaving has no column of its own --- every column being
measured \emph{from} it --- so $q_1$ is promoted by a deletion followed by a
rank-one update. In a compiled language these are Givens sweeps; in an interpreted
one they have to be the library's, or a single vectorised refactorisation will beat
them and the measurement will be of the interpreter rather than of the data
structure.

Two things keep this from being a pure win. The dependent supports
Section~\ref{sec:method} permits are exactly what an economic $QR$ cannot hold ---
an update refuses a column already in the span --- so those faces are met by
refactorising. And the repair is not cheaper at every size: on the clouds of
Section~\ref{sec:numerics} it is worth a factor of $3.5$ at $d=8\,000$, is at
parity near $d=100$, and loses below that, so an implementation dispatches on $d$
rather than committing to one route.

\section{A numerical comparison}
\label{sec:numerics}

Table~\ref{tab:highd} fixes $n=10^3$ and takes $d$ from $10^3$ to
$1.6\cdot10^4$, the regime~\cite{fischer2003} was written for and the first in
which the linear algebra rather than the combinatorics decides the cost. Two
quantities are reported, because either alone would mislead. The first is the
relative enclosure error $\max_i\norm{p_i-x}/R-1$, a deterministic property of the
returned pair, free of timing noise, where a positive value means the ball does
not in fact contain the cloud. The second is wall-clock time, which must be read
as a comparison of \emph{implementations}: both methods spend most of each
iteration in one vectorised sweep and one repair of the factorisation, so the step
counts beside them are what carry over to another language.

Before that, Table~\ref{tab:reference} says what the cone program~\eqref{eq:primal}
costs when handed to a solver directly, since a comparison against two active-set
methods is only informative if the obvious alternative is priced. Clarabel is given
the same clouds and the program assembled by hand --- $n$ second-order cones of
dimension $d+1$, no modelling layer in between.

\begin{table}[h]
\centering
\small
\setlength{\tabcolsep}{5pt}
\begin{tabular}{r rrr rrr r}
  \toprule
  & \multicolumn{3}{c}{enclosure error} & \multicolumn{3}{c}{time (s)} & \\
  \cmidrule(lr){2-4}\cmidrule(lr){5-7}
  $d$ & act.\ set & pivoting & Clarabel
      & act.\ set & pivoting & Clarabel & ratio \\
  \midrule
  $62$ & $0$ & $0$ & $-1.4\cdot10^{-11}$ & 0.003 & 0.006 & 0.5 & $165$ \\
  $125$ & $0$ & $0$ & $-9.1\cdot10^{-10}$ & 0.005 & 0.006 & 1.1 & $231$ \\
  $250$ & $0$ & $0$ & $8.4\cdot10^{-12}$ & 0.026 & 0.032 & 98.7 & $3789$ \\
  $500$ & $0$ & $0$ & $-5.8\cdot10^{-10}$ & 0.056 & 0.058 & 193.5 & $3433$ \\
  $1\,000$ & $-4.4\cdot10^{-16}$ & $-2.2\cdot10^{-16}$ & $-1.9\cdot10^{-10}$ & 0.111 & 0.167 & 359.6 & $3251$ \\
  \bottomrule
\end{tabular}
\caption{The cone-program reference at $n=10^3$: one standard normal cloud per
row, the same cloud given to all three methods. \emph{Ratio} is Clarabel's time
over the active-set method's. The two active-set methods are timed as the median
of three runs; Clarabel once per row above $d=125$, where a single solve already
runs to minutes. One run is one draw, and Clarabel's cost is the product of a
per-iteration solve that grows with $d$ and an iteration count that does not, so
these times carry the spread of that second factor and no trend should be read
off five of them.}
\label{tab:reference}
\end{table}

Two things separate the reference from the two methods this note is about, and
only one of them is speed. Clarabel is two to three orders of magnitude slower on
every row --- it factors a KKT system built from $n(d+1)$ rows, where neither
active-set method ever factorises anything of size $n$. But the errors differ in
\emph{kind}. An interior-point method stops at a tolerance, so its ball is right
to a few parts in $10^{9}$ and, at $d=250$, is very slightly wrong in the
direction that matters: the enclosure error is positive, so the returned ball does
not quite contain the cloud. Both active-set methods terminate at a vertex and are
exact to the last bit on every row. That difference is structural rather than a
matter of tuning, and it is the reason the shipped solvers are not wrappers around
a cone solver.

\begin{table}[h]
\centering
\small
\setlength{\tabcolsep}{5pt}
\begin{tabular}{r rrr rr rrr}
  \toprule
  & \multicolumn{3}{c}{steps} & \multicolumn{2}{c}{enclosure error} & \multicolumn{3}{c}{time (s)} \\
  \cmidrule(lr){2-4}\cmidrule(lr){5-6}\cmidrule(lr){7-9}
  $d$ & support & act.\ set & pivoting & act.\ set & pivoting
      & act.\ set & pivoting & ratio \\
  \midrule
  $1\,000$  &  92 &  92 & 197 & $-4.4\cdot10^{-16}$ & $0$                 & 0.148 & 0.232 & 1.56 \\
  $2\,000$  & 124 & 128 & 208 & $-1.1\cdot10^{-16}$ & $-1.1\cdot10^{-16}$ & 0.302 & 0.355 & 1.18 \\
  $4\,000$  & 177 & 177 & 249 & $-6.7\cdot10^{-16}$ & $0$                 & 1.296 & 1.450 & 1.12 \\
  $8\,000$  & 217 & 217 & 328 & $-1.0\cdot10^{-15}$ & $0$                 & 2.760 & 3.272 & 1.19 \\
  $16\,000$ & 283 & 283 & 352 & $-1.4\cdot10^{-15}$ & $0$                 & 7.869 & 8.457 & 1.07 \\
  \bottomrule
\end{tabular}
\caption{Growth in $d$ at $n=10^3$: standard normal clouds, one per row, median of
three runs. Apple M4 Pro, CPython 3.12.13, NumPy 2.5.1, SciPy 1.18.0. Both methods
repair the $QR$ of their support's edge matrix rather than rebuilding it, and both
identify the same support set, whose size is the first column. \emph{Steps} are
active-set iterations against pivot steps, which count comparable work here: one
solve or projection of order $r$ plus one $O(nd)$ sweep each. Neither reference is
run here: Welzl's median basis count is already $2.8\cdot10^5$ at $d=11$, and
Clarabel --- already minutes per solve at $d=10^3$ in Table~\ref{tab:reference} ---
would be handed a program of $n(d+1)\approx1.6\cdot10^7$ rows to factor.}
\label{tab:highd}
\end{table}

The step counts separate by about a factor of two, in the direction the mechanics
predict. The active-set count equals the final support size on four of the five
rows --- on Gaussian data the method adds points and essentially never drops one
--- while the pivoting method takes $1.2$ to $2.1$ times as many steps, the price
of never taking a full step: each walk is truncated by the first point of $S$ to
touch the shrinking sphere. The wall clocks do not separate by anything like that
factor. Their ratio sits between $1.1$ and $1.6$ with no trend in $d$, and repeat
runs move it by a few percent, so it is a small constant rather than a curve: with
$r$ in the hundreds both methods spend most of each iteration in the same two
places, and the pivoting method's extra steps are its cheapest ones. Since both
columns were measured with the factorisation repaired rather than rebuilt, the
ratio between them measures what the two searches cost and not what two data
structures cost.

Two further things are worth naming. The support grows like $\sqrt d$ --- $92$
points at $d=10^3$ and $283$ at $d=1.6\cdot10^4$ --- staying an order of magnitude
below the Carath\'eodory bound $d+1$, so the per-iteration solve is of that order
rather than of $d$; with $O(\sqrt d)$ iterations of $O(nd)$ each the total should
scale as $d^{3/2}$, and it does, a $53$-fold increase in time over a $16$-fold
increase in $d$. And the enclosure errors stay at a few ulp for both methods
across four doublings of $d$: the exactness is a property of terminating at a
basis, not of small $d$.

\section{Which to use}
\label{sec:cost}

Per iteration each method performs one $O(nd)$ sweep of the cloud --- squared
distances to the centre in Section~\ref{sec:method}, the stopping
fractions~\eqref{eq:fgk-step} in Section~\ref{sec:fgk} --- and one solve or
projection of order at most $d$. Nothing of size $n$ is ever factorised, which is
the structural difference from an interior-point method applied
to~\eqref{eq:primal}, and both solve the final face exactly rather than leaving
the caller to decide which slacks of size $10^{-9}$ were meant to be zero.

What they can hand back differs, and that is the practical reason to keep both.
The pivoting method is primal feasible at every step: interrupt it at any
iteration and what is on hand is an enclosing ball, correct if not yet smallest,
with a radius that only falls from here. That is the guarantee to want under a
time budget, and the dual method cannot offer it --- by
Proposition~\ref{prop:duality} its radius is a lower bound throughout. In exchange
the dual method finishes with the weights, which are the conic dual
of~\eqref{eq:primal} after one rescaling: $z_i=u_i(1,\,-(p_i-x)/R)$ lies in
$\mathcal{Q}^{d+1}$, is tight exactly on the support, and pairs with the slack
$s_i=(R,\,p_i-x)$ to give $s_i\T z_i=0$. A caller can verify optimality
from~\eqref{eq:kkt} in one pass over the data, without re-solving. It is also the
faster of the two on every cloud measured, which is why \texttt{cvxball} makes it
the default.

Both methods are available at \url{https://github.com/tschm/cvxball}, sharing the
factorisation of Section~\ref{sec:qr}, with the Clarabel cone program and Welzl's
recursion beside them as references.

\begin{thebibliography}{9}
\footnotesize
\setlength{\itemsep}{0pt}
\setlength{\parskip}{0pt}
\bibitem{cavaleiro2018} M.~Cavaleiro and F.~Alizadeh. A faster dual algorithm
  for the Euclidean minimum covering ball problem. \emph{Annals of Operations
  Research}, 2018. DOI 10.1007/s10479-018-3123-5.
\bibitem{cavaleiro2021} M.~Cavaleiro and F.~Alizadeh. A dual simplex-type
  algorithm for the smallest enclosing ball of balls. \emph{Computational
  Optimization and Applications} 79(3):767--787, 2021.
\bibitem{dearing2009} P.~M.~Dearing and C.~R.~Zeck. A dual algorithm for the
  minimum covering ball problem in $\R^n$. \emph{Operations Research Letters}
  37(3):171--175, 2009.
\bibitem{elzinga1972} J.~Elzinga and D.~W.~Hearn. Geometrical solutions for some
  minimax location problems. \emph{Transportation Science} 6(4):379--394, 1972.
\bibitem{fischer2003} K.~Fischer, B.~G\"artner and M.~Kutz. Fast
  smallest-enclosing-ball computation in high dimensions. \emph{Proc.\ ESA},
  630--641, 2003.
\bibitem{gaertner1999} B.~G\"artner. Fast and robust smallest enclosing balls.
  \emph{Proc.\ ESA}, 325--338, 1999.
\bibitem{gaertner2000} B.~G\"artner and S.~Sch\"onherr. An efficient, exact,
  and generic quadratic programming solver for geometric optimization.
  \emph{Proc.\ 16th Annu.\ ACM Sympos.\ Comput.\ Geom.}, 110--118, 2000.
\bibitem{sylvester1857} J.~J.~Sylvester. A question in the geometry of
  situation. \emph{Quarterly Journal of Pure and Applied Mathematics} 1:79, 1857.
\bibitem{welzl1991} E.~Welzl. Smallest enclosing disks (balls and ellipsoids).
  \emph{New Results and New Trends in Computer Science}, LNCS 555, 359--370, 1991.
\end{thebibliography}

\end{document}
